\section{Úvod}

V modernom vývoji softvéru sa architektúra mikroservisov stala štandardom pre budovanie škálovateľných a flexibilných systémov~\cite{Newman2021}. Na rozdiel od monolitických aplikácií, kde komunikácia prebieha v rámci jedného pamäťového priestoru, mikroservisy intenzívne využívajú počítačovú sieť. Sieť však z hľadiska distribuovaných systémov neponúka absolútnu spoľabilivosť — bežne dochádza ku kolísaniu latencie, strate paketov alebo dočasným výpadkom uzlov.

Pokiaľ nie je komunikačná vrstva správne navrhnutá, lokálna sieťová porucha môže viesť ku kaskádovému zlyhaniu celého systému. Táto práca analyzuje správanie rôznych transportných protokolov a spracovateľov správ (message brokers) v podmienkach sieťovej degradácie a vyhodnocuje vplyv architektúrnych vzorov odolnosti (Resilience Patterns) na zachovanie dostupnosti aplikácie.

Základný problém a motivácia výskumu sú detailnejšie popísané v časti~\ref{s: problem}.
Navrhnutá architektúra testovacieho stojanu a metodika Chaos Engineeringu sú predstavené v časti~\ref{s: architektura}.
Porovnanie komunikačných protokolov a vzorov odolnosti prinášajú časti~\ref{s: protokoly} a~\ref{s: vzory}.
Záverečné zhrnutie a inžinierske odporúčania obsahuje časť~\ref{s: zaver}.